\section{池化层}

\subsection{池化的基本概念}

\begin{figure}[H]
\centering
\includegraphics[width=0.95\textwidth]{slides-images/slide-057.jpg}
\caption{池化层：对每个通道独立进行空间下采样\protect\footnotemark}
\end{figure}
\footnotetext{Slides 第58页。}

池化层是 CNN 中另一种实现\textbf{空间下采样}的方式。与带步幅的卷积不同，池化层：
\begin{itemize}
    \item 没有可学习的参数
    \item 对每个通道\textbf{独立}操作
    \item 保持通道数不变，只缩小空间尺寸
\end{itemize}

\subsection{最大池化（Max Pooling）}

\begin{figure}[H]
\centering
\includegraphics[width=0.85\textwidth]{slides-images/slide-058.jpg}
\caption{$2 \times 2$ 最大池化（步幅 2）：在每个不重叠的 $2 \times 2$ 区域内取最大值\protect\footnotemark}
\end{figure}
\footnotetext{Slides 第59页。}

最常见的池化操作是 $2 \times 2$ 最大池化，步幅为 2：
\begin{itemize}
    \item 将每个通道的特征图划分为不重叠的 $2 \times 2$ 小块
    \item 在每个小块内取\textbf{最大值}
    \item 空间尺寸缩小为原来的一半
\end{itemize}

其他选择包括平均池化（Average Pooling）和抗锯齿下采样（Anti-alias Downsampling）。

\begin{knowledgebox}{池化层的实践建议}
\begin{itemize}
    \item 最常用的配置：$2 \times 2$ kernel，stride 2（``我想把空间尺寸减半''）
    \item 池化层通常\textbf{不使用填充}——对于最大池化，填充零等价于 ReLU，是多余的
    \item 现代架构越来越倾向于用\textbf{带步幅的卷积}代替池化，因为前者有可学习参数
\end{itemize}
\end{knowledgebox}

\subsection{全局平均池化}

在 CNN 的末尾，常用\textbf{全局平均池化}（Global Average Pooling）将每个通道的整个空间维度压缩为一个数值。如果最后的特征图大小为 $C \times H \times W$，全局平均池化输出 $C \times 1 \times 1$，直接用于分类。

\subsection{本章小结}

池化层通过无参数的空间下采样操作缩小特征图尺寸，从而增大后续层的感受野并减少计算量。最大池化是最经典的选择，但现代架构中带步幅的卷积越来越多地取代了池化层的角色。

%% ========== Part 7: 构建完整的 CNN ==========

